\section{Causal affordances and executable tests for an agent-assisted library}
\label{sec:exec24-affordance-fuzzing}
A proposed causal assistant must select useful analytical actions and expose mistakes in the software that carries them out. These are related but distinct responsibilities. Reasoning about which familiar object can serve an unfamiliar purpose offers a controlled test of functional selection; generating programs that expose computational inconsistencies offers a test of software behavior. Neither alone establishes a banking intervention effect. The affordance study by Gjerde, Cheung and Lagnado and the TitanFuzz study by Deng and colleagues motivate a joint design in which agents propose actions and test inputs while explicit scientific and computational contracts determine what counts as success.\citep{exec24_affordance_original,exec24_fuzzing_original,exec24_affordance_fuzzing_reference}

\subsection{What an affordance benchmark measures}
The inspected affordance text is the 21-page February 2025 author preprint. It studies 20 tasks in which an ordinary tool is unavailable and the respondent selects a substitute. The four-choice design includes objects associated with the ordinary tool but unsuitable for the task, a less associated suitable object and a distractor. Experiment one recruits 100 UK participants and compares their responses with historical model configurations under different prompts and temperatures. The paper explicitly acknowledges that repeated stochastic model runs are not independent human participants, although run-level averages are used as comparable data points.\citep{exec24_affordance_original}

This is a useful challenge to association-based selection: the familiar object can be the wrong answer. It does not identify an external intervention effect from observational data. Task authors have specified the intended function and scored the offered alternatives; a bank still needs evidence about legal eligibility, timing, collateral, assignment and outcome mechanisms before translating a plausible suggestion into an action. The study's qualitative examples illuminate response patterns, but they were selected for notable differences and are not a separate representative sample of reasoning failures.\citep{exec24_affordance_original,exec24_affordance_fuzzing_reference}

Experiment two retains 296 of 300 recruited UK participants and introduces nine-option distractor and image conditions. Its ordinary prompt permits longer answers than the first experiment's restricted prompt. The authors explicitly discuss this change as a possible explanation of part of the earlier prompting effect. The separate-option reasoning prompt is not evaluated in the nine-option condition because the resulting answers would be long. This design does not support a fully crossed estimate of every prompt-by-difficulty interaction.\citep{exec24_affordance_original}

The reported results include configuration-specific similarities to human accuracy, performance declines with additional distractors, and different responses to images across models. A nonsignificant difference is not an equivalence result: equivalence requires a justified margin and an appropriate test. The paper also distinguishes newly written tasks from structurally new problems and acknowledges incomplete control over training-data overlap. Its behavioral observations do not determine the internal mechanism or demonstrate real-world tool use. Historical model names describe the studied conditions, not current product specifications or a claim of superiority today.\citep{exec24_affordance_original,exec24_affordance_fuzzing_reference}

\subsection{Translate functional selection into a banking contract}
Consider a stipulated liquidity problem: obtain eligible funds within one business day. An action's required properties are a settlement delay no greater than one day, eligible collateral and an existing authorization. Four authored choices are a familiar facility settling in two days, an unfamiliar facility satisfying all three conditions, an immediate facility without eligible collateral, and an immediate facility lacking authorization. The unfamiliar facility is the unique feasible choice. Adding five infeasible distractors leaves that answer unchanged. These are declared fixture properties, not claims about actual facilities or law.\citep{exec24_affordance_fuzzing_reference}

The assistant can decompose the request into constraints and attach source passages, dates and unresolved interpretations to each property. It should then distinguish feasibility from preference. If two choices satisfy the constraints, ranking them requires costs, risk and the bank's objective; if collateral eligibility is unknown, the answer must retain that uncertainty. A familiar label or an explanation that sounds financially plausible cannot supply the missing fact. In graph construction, the analogous step is to distinguish candidate variables and functional mechanisms from assumptions needed for causal identification.\citep{exec24_affordance_fuzzing_reference}

Test representation and substantive meaning separately. Renaming facilities or permuting choices should preserve this fixture's answer; changing settlement from one to two days should not. A unit conversion must preserve the same economic threshold. Changing an observation timestamp can invalidate a proposed covariate. These transformations create an auditable test family whose required behavior is defined before a model generates an answer. They do not establish that the chosen public task distribution represents a bank's future decisions.\citep{exec24_affordance_fuzzing_reference}

Choice count also changes a naive accuracy baseline. In an authored comparison, accuracy 0.85 among four choices has chance-normalized score
\[
 (0.85-1/4)/(1-1/4)=0.80,
\]
whereas accuracy 0.75 among nine choices gives
\[
 (0.75-1/9)/(1-1/9)=23/32=0.71875.
\]
Both raw and normalized scores are informative; neither removes changes in task difficulty or establishes causal competence. The numbers are illustrative and are not the paper's reported contrast.\citep{exec24_affordance_fuzzing_reference}

Repeated outputs do not automatically increase the number of independent problems. Suppose there are 20 independent task clusters, each with 25 unit-variance errors, equicorrelated at a stipulated 0.60 within task. The mean over 500 outputs has variance
\[
 \frac{1+(25-1)0.60}{20\cdot25}=0.0308,
\]
compared with 0.002 under independence. The equivalent independent count for this variance calculation is $2500/77$, about 32.47. A different dependence law gives a different answer. This witness motivates retaining task identity, shared context, prompt selection and model revision in the evaluation design. It is not an estimate of the study's correlation or a reanalysis of its significance tests.\citep{exec24_affordance_fuzzing_reference}

\subsection{Generate difficult programs while preserving the test oracle}
TitanFuzz supplies a complementary engineering idea. Its inspected March 2023 author version combines a generative model for seed programs with an infilling model for mutations. Mutation operators alter arguments, surrounding code and methods; adaptive operator selection rewards valid generations, while a structural fitness score favors dataflow depth and unique API calls and penalizes repetition. The paper uses successful execution with target invocation as a program-validity criterion. Validity, diverse execution and discovery of a real defect are separate outcomes.\citep{exec24_fuzzing_original}

The paper's principal oracles compare CPU and GPU computations and identify crashes. Comparison tolerances address ordinary numerical variation. The experiments cover historical PyTorch and TensorFlow configurations, with a one-minute per-API budget for the main comparison and a four-minute budget for bug finding. Reported line coverage concerns the measured Python code; it does not establish that every lower-level path, relevant banking contract or possible defect was exercised. The paper reports 65 candidate findings, 53 developer confirmations and 41 previously unknown bugs at its reporting date, with eight rejected findings. One rejected example has intentionally implementation-dependent behavior. This distinction matters: a discrepancy is evidence to investigate, not automatically a valid defect.\citep{exec24_fuzzing_original}

The inspected repository revision is dated September 2023 and reports later totals of 55 confirmations and 44 previously unknown bugs. These dated source statements remain separate; we have not independently adjudicated the individual reports. The revision's complete tree has 88 entries. Its README identifies the research paper, describes supplied historical seed artifacts and points to a seed-generation entrypoint absent from that tree. The external seed archive has not been acquired in this review. Thus the repository does not by itself establish complete reproduction of the paper's generation process.\citep{exec24_titanfuzz_source}

A version contract is also necessary. The inspected requirements leave several major dependencies unpinned, while the Docker recipe names TensorFlow 2.11.0rc1, differing from the paper's TensorFlow 2.10 experiment. The source includes process recovery, timeout and skip statuses, but these mechanisms alone do not establish containment, reproducible numerical behavior or target-platform acceptance. No native model, downloaded program, container or package is executed in this review.\citep{exec24_titanfuzz_source}

\subsection{A numerical agreement check is not an economic acceptance rule}
The inspected comparator declares relative tolerance 0.01 and absolute tolerance 0.001 and calls library closeness functions with equal-NaN handling. Some unhandled value types fall through to a true return. These are static source observations whose consequences need carefully specified native tests. They should not be converted into a claim that a particular version has been observed to misclassify a banking result.\citep{exec24_titanfuzz_source}

An authored scalar example illustrates the required distinction. Define the comparison rule
\[
 |x-y|\leq a+r|y|,
\]
with $a=0.001$ and $r=0.01$. If two backends return 1,000,000 USD and 1,005,000 USD, their difference is 5,000 USD and the tolerance is 10,050.001 USD. This rule accepts agreement even though the difference exceeds a separately stipulated 1,000 USD business tolerance. The scalar rule is our explicit witness using the inspected constants; it is not a native execution of either numerical library.\citep{exec24_affordance_fuzzing_reference}

Absolute tolerances carry units. Comparing zero with 0.0005 dollars passes an absolute threshold of 0.001 dollars. Expressing the same values in cents while retaining the numerical threshold 0.001 fails; the threshold must also be converted. A causal API therefore needs units, estimand and error budgets in its artifact contract, including any conversion performed by a CLI or MCP consumer.\citep{exec24_affordance_fuzzing_reference}

Backend agreement can also preserve a shared error. If the exact truth is 100, both backends return 110 and action occurs above 105, perfect agreement produces the wrong action. In inference, a shared identification mistake is even less accessible to numerical comparison. Use independent mathematical references, exact finite models, invariance tests and deliberately invalid requests alongside differential tests. Refusal on an unidentified effect is part of correct behavior, not a program-generation failure to discard.\citep{exec24_affordance_fuzzing_reference}

\subsection{Adoption, interfaces and the public prototype}
The defensible proposal is to adopt the testing pattern conditionally while treating the available code as a separately qualified dependency. No license, copying or notice file appears in the inspected tree, and the conventional license request returns 404. This bounded observation leaves reuse rights unresolved; it does not establish that no permission exists elsewhere. Direct code adoption or porting should wait for an applicable grant and review of dependencies, model artifacts and supplied seeds. A project-owned harness can be designed from the inspected scientific ideas without declaring the downloaded implementation accepted.\citep{exec24_titanfuzz_source,exec24_affordance_fuzzing_reference}

A public prototype should freeze a versioned task schema before testing: business objective, units, clocks, graph and identification contract, permitted operations, expected success or refusal, and the numerical tolerance. The agent proposes candidate programs or structured requests. A controlled runner records their source identity, environment, resource limits and statuses; an independently specified oracle evaluates outputs. The scientific checker must be allowed to reject a numerically successful calculation when the causal target is unsupported. Corpus revisions and model changes require a new held-out comparison.\citep{exec24_affordance_fuzzing_reference}

For CLI and MCP consumers, generation, execution, adjudication and reporting should exchange separate immutable artifacts. A request identifies its oracle and authority scope; a response records all attempts, timeouts, duplicates, failures, unresolved discrepancies and accepted outputs. Sibling repositories are prospective consumers of this protocol, not demonstrated integrations. Before copying a discovered reproducer into another workspace, review its permissions, inputs, resource behavior and dependence on the environment. These are proposed product obligations, with no cross-repository deployment asserted.\citep{exec24_affordance_fuzzing_reference}

An exact cost fixture starts with 100 attempts, 60 runnable programs and 20 duplicates among the runnable outputs, leaving 40 unique runnable programs. Ten candidate discrepancies comprise six confirmed defects, two intended differences and two unresolved cases. At a stipulated total cost of 600 USD, costs are 6 USD per attempt, 10 USD per runnable program, 15 USD per unique runnable program and 100 USD per confirmed defect. These denominators include different services. They are not market prices or a measured return on investment. A future comparison must include generation, hardware, containment, failed runs and human adjudication, then connect accepted outputs to decision value.\citep{exec24_affordance_fuzzing_reference}

The combined lesson is practical. Functional decomposition can help formulate an analysis and adversarial program generation can expose implementation errors. The merit of a public experiment comes from its explicit semantics, independent oracle and complete accounting. Neither task accuracy nor software coverage supplies missing assignment information, resolves observationally equivalent causal worlds or substitutes for later institutional acceptance.\citep{exec24_affordance_fuzzing_reference}
